\section{\label{sec:spdf}Relation to light-front distributions}

We discuss the relation between quasi and light-front PDFs by examining the 
Mellin moments of these distributions, and using the connection between 
Mellin moments and matrix elements of local operators, which are twist-two in 
the case of light-front PDFs \cite{Christ:1972ms}. For the quasi PDFs, the 
local operators corresponding to the Mellin moments do not have a well-defined 
twist, but can be related to twist-two operators after subtracting higher 
twist effects and applying target-mass 
corrections \cite{Lin:2014zya,Alexandrou:2015rja}. Although we 
consider smeared matrix elements in this work, the arguments regarding higher twist and 
target mass effects in \cite{Lin:2014zya,Alexandrou:2015rja} still apply, 
because the flow time serves as an alternative
gauge-invariant regulator to the lattice spacing.

We connect the Mellin moments of the quasi PDF to matrix elements of 
local operators in the following way. Working in axial gauge, $B_z(x;\tau)=0$, 
the 
matrix element $h^{(s)}$ is
\begin{equation}\label{eq:hAzdef}
h^{(s)}\left(\frac{z}{\sqrt{\tau}},\sqrt{\tau}P_z,\sqrt{\tau}\Lambda_{\mathrm{
QCD}},
\sqrt{\tau}
M_\mathrm{N}\right)\Big|_{B_z=0} = 
\frac{1}{2P_z}\left\langle 
P_z \left| \overline{\chi}(z;\tau) 
\gamma_z\frac{\lambda^a}{2}\chi(0;\tau)\right| 
P_z\right\rangle_\mathrm{C}.
\end{equation}
We now substitute this expression into the definition of the quasi PDF, 
Equation \eqref{eq:qpdfdef}, and integrate the resulting
expression over the full range of $\xi$. In contrast to the light-front 
PDF, this range extends from negative to positive infinity, giving
\begin{equation}
\int_{-\infty}^\infty 
\mathrm{d}\xi\, 
q^{\,(s)}\left(\xi,\sqrt{\tau}P_z,\sqrt{\tau}\Lambda_{\mathrm{QCD}},\sqrt{\tau}
M_\mathrm{N}\right)\Big|_{B_z=0} 
= 
h^{(s)}(0,\sqrt{\tau}P_z,\sqrt{\tau}\Lambda_{\mathrm{QCD}},\sqrt{\tau}M_\mathrm{
N}
)\Big|_{B_z=0}.
\end{equation}
Here we have used the relation $\delta(zP_z) = 
\delta(z)/P_z$, for $P_z>0$. We see that the first Mellin moment of the quasi 
PDF can be expressed in terms of the Euclidean matrix element of a local 
(smeared) operator.

We extend this argument to arbitrary 
moments of quasi PDFs by considering derivatives of the quasi distribution with 
respect to the spatial separation $z$ \cite{Collins:2011zzd}. Inverting the 
Fourier transform in Equation \eqref{eq:qpdfdef}, we have
\begin{equation}\label{eq:handqpdf}
h^{(s)}\left(\frac{z}{\sqrt{\tau}},\sqrt{\tau}P_z,\sqrt{\tau}\Lambda_{\mathrm{
QCD}},
\sqrt{\tau}
M_\mathrm{N}\right)
=
\int_{-\infty}^\infty \mathrm{d}\xi\, e^{-i\xi z 
P_z} 
q^{\,(s)}\left(\xi,\sqrt{\tau}P_z,\sqrt{\tau}\Lambda_{\mathrm{QCD}},\sqrt{\tau}
M_\mathrm{N}\right).
\end{equation}
Applying derivatives with respect to the displacement $z$, we obtain
\begin{align}\label{eq:hint1}
\left(\frac{i}{P_z}\frac{\partial}{\partial 
z}\right)^{n-1} 
h^{(s)}{} & 
\left(\frac{z}{\sqrt{\tau}},\sqrt{\tau}P_z,\sqrt{\tau}\Lambda_{\mathrm{QCD}},
\sqrt{
\tau }
M_\mathrm{N}\right) = \nonumber\\
{} & \int_{-\infty}^\infty \mathrm{d}\xi\, \xi^{n-1}e^{-i\xi z 
P_z} 
q^{\,(s)}\left(\xi,\sqrt{\tau}P_z,\sqrt{\tau}\Lambda_{\mathrm{QCD}},\sqrt{\tau}
M_\mathrm{N}\right).
\end{align}

Defining the moments of the smeared quasi PDF, $b_n^{(s)}$, as
\begin{equation}\label{eq:bdef}
b_n^{(s)}\left(\sqrt{\tau}P_z,\frac{\Lambda_{\mathrm{QCD}}}{P_z}, 
\frac{M_\mathrm{N}}{P_z}\right) 
= \int_{-\infty}^\infty 
\mathrm{d}\xi\, \xi^{n-1} 
q^{\,(s)}\left(\xi,\sqrt{\tau}P_z,\sqrt{\tau}\Lambda_{\mathrm{QCD}},\sqrt{\tau}
M_\mathrm{N}\right),
\end{equation}
and substituting the definition of the matrix element $h^{(s)}$, given in 
Equation \eqref{eq:hdef}, into Equation \eqref{eq:hint1}, in the limit that 
$z\rightarrow 0$ we obtain
\begin{align}
b_n^{(s)}\left( \sqrt{\tau}P_z,\frac{\Lambda_{\mathrm{QCD}}}{P_z}, 
\frac{M_\mathrm{N}}{P_z}\right)_{B_z=0}
= {} & \frac{c_n^{(s)}(\sqrt{\tau}P_z)}{2P_z^n} \nonumber\\
{} & \qquad \times \left\langle 
P_z \left| \left[\overline{\chi}(z;\tau) \gamma_z 
\left(i\overleftarrow{\partial}_z^{n-1}\right)\frac{\lambda^a}{2}
\chi(0;\tau)\right ] _{z=0 } \right|P_z\right\rangle_\mathrm{C}.
\end{align}
The perturbative coefficients, $c_n^{(s)}(\sqrt{\tau}P_z)$, capture potential 
singularities in the righthand side of Equation \eqref{eq:hint1} in the limit 
of vanishing separation z and vanishing flow time $\tau$, and follow from a smeared operator product expansion 
\cite{Monahan:2015lha} approach to the nonlocal matrix element, as outlined in 
\cite{Lin:2014zya}.

We restore gauge invariance to obtain our final expression for the moments of 
quasi PDFs in terms of Euclidean matrix elements of local operators:
\begin{equation}
b_n^{(s)} \left(\sqrt{\tau}P_z,\frac{\Lambda_{\mathrm{QCD}}}{P_z}, 
\frac{M_\mathrm{N}}{P_z}\right) 
=\frac{c_n^{(s)}(\sqrt{\tau}P_z)}{2P_z^n}\left\langle 
P_z \left| \left[
\overline{\chi}(z;\tau)\gamma_z (i\overleftarrow{D}_z)^{(n-1)} 
\frac{\lambda^a}{2}\chi(0;\tau) \right ] _{z=0 }
\right|P_z\right\rangle_\mathrm{C} .
\end{equation}

The local operators that appear in the matrix element on the right hand side of 
this expression are not twist-two operators: they are not symmetric and 
traceless. The discrepancy between these matrix elements and matrix elements of 
twist-two operators are given by corrections  that appear in powers of 
$\Lambda_\mathrm{QCD}^2/{P_z^2}$ 
and $M_\mathrm{N}^2/P_z^2$ \cite{Lin:2014zya,Alexandrou:2015rja}. The terms 
that scale as ${\cal O}(M_\mathrm{N}^2/P_z^2)$ correspond to target mass 
corrections \cite{Nachtmann:1973ll,Georgi:1974sr}.  Although the appropriate 
interpretation 
of PDFs in the presence of these target mass corrections is subtle 
\cite{Steffens:2012jx,Monfared:2014nta}, for our purposes it is sufficient that 
these 
non-leading 
corrections can be absorbed 
by writing \cite{Lin:2014zya,Alexandrou:2015rja}
\begin{align}
b_n^{(s)} \left(\sqrt{\tau}P_z,\frac{\Lambda_{\mathrm{QCD}}}{P_z}\right) 
= {} & \frac{c_n^{(s)}(\sqrt{\tau}P_z)}{2P_z^n}\left\langle 
P_z \left|\left[
\overline{\chi}(z;\tau)\gamma_z (i\overleftarrow{D}_z)^{(n-1)} 
\frac{\lambda^a}{2}\chi(0;\tau) \right ] _{z=0 }
\right|P_z\right\rangle_\mathrm{C} \nonumber \\
{} & \qquad \times K_n^{-1}\left(\frac{M_\mathrm{N}^2}{4P_z^2}\right),
\end{align}
where
\begin{equation}
K_n\left(\frac{M_\mathrm{N}^2}{4P_z^2}\right) = \sum_{j=0}^{n/2} 
\bigg(
\begin{array}{c}
n-j \\
j
\end{array}
\bigg)
\left(\frac{M_\mathrm{N}^2}{4P_z^2}\right)^j.
\end{equation}

The corrected matrix 
elements on the right hand side of this equation can now 
be expanded in a Taylor series with respect to
$\Lambda_\mathrm{QCD}^2/P_z^2$. The coefficients in this expansion 
represent higher twist effects that arise because the original matrix 
element is not a matrix element of a twist-two operator. The leading 
coefficient 
in this expansion is a twist-two contribution that can depend only on the 
nucleon structure and the flow time:
\begin{equation}\label{eq:btwist}
b_n^{(s)} \left(\sqrt{\tau}P_z,\sqrt{\tau}\Lambda_{\mathrm{QCD}}\right)  
= 
c^{(s)}_n(\sqrt{\tau}P_z)b_n^{(s,\mathrm{twist}-2)}
\left(\sqrt{\tau}\Lambda_{\mathrm{
QCD}}\right)  +{\cal 
O}\left(\frac{\Lambda_\mathrm{QCD}^2}{P_z^2}\right),
\end{equation} 
where, for $\Lambda_\mathrm{QCD}^2/P_z^2\ll 1$, the higher twist corrections 
can be ignored.

In summary, we assume that: first, we can correct exactly for target mass 
corrections; and second, we can take the momentum $P_z$ sufficiently large that 
higher 
twist effects are negligible. Then, under these 
assumptions, the moments of the smeared quasi PDFs are dimensionless products 
of perturbative coefficients and pure twist-two matrix elements, which are only 
functions of the 
dimensionless quantity
$\sqrt{\tau} \Lambda_{\mathrm{QCD}}$, that contain
information about the structure of the hadron.

\subsection{\label{sec:renormpdf} Short-distance expansion}

We can now relate the moments of the smeared quasi PDF 
$b_n^{(s,\mathrm{twist}-2)}\big(\sqrt{\tau}\Lambda_{\mathrm{QCD}}\big)$, which 
are local matrix elements of smeared fields, to the renormalized moments of the 
light-front PDFs, by using the properties of the gradient flow that arise from 
a short distance expansion 
\cite{Luscher:2011bx,Suzuki:2013gza,Asakawa:2013laa,Makino:2014taa,
Hieda:2016lly}. 
The exponentially local nature of the smearing procedure 
allows for a short distance expansion of the smeared local operators in terms 
of renormalized operators in some renormalization scheme, such as the 
$\overline{MS}$ scheme. It is straightforward to show that this short 
distance expansion leads to 
\begin{equation}\label{eq:shortdist}
b_n^{(s,\mathrm{twist}-2)}\big(\sqrt{\tau}\Lambda_{\mathrm{QCD}}\big) = 
\widetilde{C}_n^{(0)}(\sqrt{\tau}\mu)a^{(n)}(\mu) + {\cal 
O}(\sqrt{\tau}\Lambda_{\mathrm{QCD}}),
\end{equation}
where $\mu$ is a renormalization scale. The leading order term in this 
expansion, $a^{(n)}(\mu)$, is the matrix element of a renormalized twist-two 
operator with the same gamma matrix and derivative structure as the smeared 
operator that appears in the matrix element on the left hand side. The higher 
order terms arise from higher dimension operators that enter the short distance 
expansion of the smeared matrix element.

We now combine this short-distance expansion with Equation \eqref{eq:btwist} to 
write
\begin{equation}\label{eq:double}
b_n^{(s)}\big(\sqrt{\tau}\Lambda_{\mathrm{QCD}}\big) = 
C_n^{(0)}(\sqrt{\tau}\mu, \sqrt{\tau}P_z)a^{(n)}(\mu) + {\cal 
O}\left(\sqrt{\tau}\Lambda_{\mathrm{QCD}},\frac{\Lambda_\mathrm{QCD}^2}{P_z^2} 
\right),
\end{equation}
Both the leading short distance coefficient function,
$ C_n^{(0)}(\sqrt{\tau}\mu, \sqrt{\tau}P_z)$, and the higher order corrections 
can be 
computed in perturbation theory, so that this approximation can be 
systematically improved. 

For the rest of this discussion, we will assume that we work in a regime in 
which there is a hierarchy of scales given by
\begin{equation}\label{eq:scales}
\Lambda_{\mathrm{QCD}},M_N\ll P_z\ll  \tau^{-1/2},
\end{equation}
so that power corrections and higher-twist effects can be ignored. We also 
assume that target mass corrections have been applied.

To relate the smeared quasi PDF with the light-front 
PDF, we introduce a kernel function, $Z(x,\sqrt{\tau}\mu,\sqrt{\tau}P_z)$,
whose Mellin moments are given by
\begin{equation}\label{eq:cinvz}
 \left[C_n^{(0)}(\sqrt{\tau}\mu, \sqrt{\tau}P_z)\right]^{-1} = 
\int_{-\infty}^{\infty} dx\, x^{n-1}  
Z(x,\sqrt{\tau}\mu, \sqrt{\tau}P_z).
\end{equation}
With this definition, and using the properties of multiplicative convolution, 
we find 
\begin{equation}
f(x,\mu) =  \int_{-\infty}^{\infty} \frac{d\xi}{\xi}\, 
Z\left(\frac{x}{\xi},\sqrt{\tau}\mu, \sqrt{\tau}P_z\right) 
q^{\,(s)}\left(\xi,\sqrt{\tau} \Lambda_{\mathrm{QCD}}\right) + {\cal 
O}(\sqrt{\tau}\Lambda_{\mathrm{QCD}}).
\end{equation}
We introduce the inverse kernel through
\begin{equation}
C_n^{(0)}(\sqrt{\tau}\mu, \sqrt{\tau}P_z) = \int_{-\infty}^{\infty} dx\, 
x^{n-1}  
\widetilde{Z}(x,\sqrt{\tau}\mu, \sqrt{\tau}P_z),
\end{equation}
which leads to
\begin{equation}
q^{\,(s)}\left(x,\sqrt{\tau} \Lambda_{\mathrm{QCD}}, \sqrt{\tau}P_z\right) =  
\int_{-1}^{1} 
\frac{d\xi}{\xi}\, \widetilde{Z}\left(\frac{x}{\xi},\sqrt{\tau}\mu, 
\sqrt{\tau}P_z\right) 
f(\xi,\mu) + 
{\cal O}(\sqrt{\tau}\Lambda_{\mathrm{QCD}})
\end{equation}
Note that all of these relations are only 
valid if
\begin{equation}\label{eq:scales_tau}
\Lambda_{QCD},M_N\ll P_z \ll \tau^{-1/2}.
\end{equation}

The kernel function can be 
computed in continuum perturbation theory, following the methods introduced 
in~\cite{Luscher:2011bx} and the examples 
in~\cite{Ma:2014jla,Ji:2015jwa,Chen:2016fxx,Ishikawa:2016znu}.
Given that those computations are performed in the continuum, 
they are independent of the lattice formulation used to extract the 
smeared quasi-PDFs introduced in this paper and lattice computations of 
smeared quasi-PDFs can be performed with a variety of lattice actions, each of 
which results in the same universal continuum quasi-PDF. Consequently, lattice 
computations are disconnected from the matching procedure of these universal, 
continuum quasi-PDFs to the light-front PDFs.  

In both our approach, as in Ji's approach, 
the large nucleon momentum serves only to suppress higher 
twist contributions. We have, however, introduced a new scale, the (inverse) 
flow time, $\tau^{-1}$,  that serves as a regulator of ultraviolet divergences 
that arise from external composite operators  that define the matrix element. The 
scale $\tau^{-1}$ needs to be large relative to hadronic scales but remains 
finite. These 
requirements for the hierarchy of scales, expressed in Equation 
\eqref{eq:scales_tau}, are no different in nature than the requirements used to 
factor physical cross-sections into PDFs and Wilson coefficients and are 
similar in spirit to the factorization approach proposed in 
\cite{Ma:2014jla,Ishikawa:2016znu}. In this approach, the renormalization 
scale, which plays a similar role as our flow time scale,
and the factorization scale are distinct and separate from the large momentum, 
which suppresses higher twist effects. 

